\chapter{ניתוח שוק}\label{ch:market}

שוק הירקות הטריים בישראל גדול, יציב ותלוי בייבוא ובשטח חקלאי מצטמצם. גגון פועלת בנישה צרה ומוגדרת היטב: ירק עלים ועשבי תבשיל טריים, שגדלים במרחק מאות מטרים מהצרכן. פרק זה מכמת את ההזדמנות, מגדיר את מגזרי היעד ומשווה את גגון לחלופות הקיימות.

\section{גודל השוק: \en{TAM}, \en{SAM}, \en{SOM}}\label{sec:market-size}

אנו מודדים את השוק בשלוש שכבות. כל הנתונים בסעיף זה הם \textbf{הערכת צוות} ולא נתוני צד שלישי. הם מבוססים על ספירת בניינים, על מחיר ממוצע של \RRprice~\shekel{} לק"ג ועל התפוקה של גג יחיד. את ההנחות יש לבחון בדיליגנס, ואנו מציגים אותן במלואן כדי שניתן יהיה לערער עליהן.

\begin{itemize}
  \item \textbf{\en{TAM}, שוק כולל}: צריכה שנתית של ירק עלים ועשבי תבשיל טריים במשק הבית ובמסעדנות הישראלית, בערך של \RRtam{} מיליון \shekel{} בשנה (הערכת צוות).
  \item \textbf{\en{SAM}, שוק בר-שירות}: החלק שמשרת מודל גגות עירוני בגוש דן. הוא כולל בניינים בני שש קומות ומעלה, מסעדות ובתי קפה ברדיוס של קילומטר אחד, ומשקי בית שמוכנים לשלם פרמיה על טריות. גודלו \RRsam{} מיליון \shekel{} בשנה (הערכת צוות).
  \item \textbf{\en{SOM}, שוק בר-השגה}: ההכנסה שהתוכנית מכוונת אליה ב-2031, \RRsom{} מיליון \shekel{}, על בסיס \RRroofsFinal{} גגות. זה מהווה כ-\RRsomShareSam\% מה-\en{SAM} (ה-\en{SOM} כולל גם הכנסות רישוי \en{RootOS}).
\end{itemize}

איור~\ref{fig:market} מציג את שלוש השכבות. הנקודה החשובה למשקיע: התוכנית אינה דורשת נתח שוק גבוה. היא דורשת ביצוע טוב של \RRroofsFinal{} גגות, כשהביקוש המקומי לכל גג הוא ביקוש צפוי שיאומת בפיילוט.

\begin{figure}[htbp]
  \centering
  \includegraphics[width=0.92\linewidth]{fig_market}
  \caption{\en{TAM}, \en{SAM} ו-\en{SOM} לשוק ירק העלים הטרי (מיליוני \shekel{} בשנה, הערכת צוות).}
  \label{fig:market}
\end{figure}

\begin{insightbox}[תובנה]
ה-\en{SOM} של \RRsom{} מיליון \shekel{} מהווה כ-\RRsomShareSam\% מה-\en{SAM}. הסיכון בתוכנית הוא ביצוע ולא גודל שוק.
\end{insightbox}

\section{מגזרי יעד}\label{sec:segments}

גגון מוכרת לשלושה מגזרים. לכל אחד תפקיד אחר במודל: מסעדות קונות נפח קבוע, תושבים קונים נוחות ושייכות, וחוות חיצוניות קונות תוכנה.

\begin{enumerate}
  \item \textbf{מסעדות ובתי קפה ברדיוס של קילומטר}. הם קונים ירק עלים ועשבי תבשיל בכמות קבועה, ומעריכים טריות ושמות של מקור. הם מהווים את עיקר הנפח ואת העוגן להכנסה צפויה.
  \item \textbf{תושבי הבניין ושכניו (\enquote{סל ירוק})}. משקי בית בעלי הכנסה בינונית-גבוהה שרוצים ירק טרי, מקומי ושקוף. הם קונים מנוי חודשי ומקבלים סל שבועי. זהו ערוץ בעל שולי רווח גבוהים ונאמנות גבוהה.
  \item \textbf{חוות ומגדלים חיצוניים (מ-2029)}. חוות קטנות ובינוניות ברישיון תוכנת \en{RootOS}. הן מהוות הכנסה חוזרת ללא הון קבוע, והן מגיעות ל-\RRsaasFarmsFinal{} חוות עד 2031.
\end{enumerate}

מגזר רביעי, ועדי בתים, אינו לקוח על המוצר אלא שותף: הוא מספק את הגג ומקבל \RRrentShare\% מהכנסות הגג. פרק~\ref{ch:business-model} מתאר את הסכם השותפות.

\section{חלופות ותחרות}\label{sec:competition}

לא זיהינו מתחרה ישיר (הערכת צוות) שמפעיל חקלאות על גגות מגורים בגוש דן בהסכם ארוך עם ועדי בתים. התחרות היא על ארנק הלקוח, מול ארבע קטגוריות של חלופות. אנו מציגים קטגוריות כלליות ואיננו מנקבים בשמות של חברות.

חוות אנכיות פנימיות (\en{indoor vertical farms}) מציעות שליטה מלאה בסביבה, אך צורכות חשמל רב לתאורה, כפי שמוצג בספרות על גידול בסביבה מבוקרת \cite{kozai2020}. חממות מסחריות זולות לשטח, אך רחוקות מהעיר. תוצרת מיובאת או שדה קונבנציונלית זולה, אך אינה טרייה ביום הקטיף. ערכות לבלקון מתאימות לעניין אישי, אך אינן מתרחבות לצריכה אמיתית. טבלה~\ref{tab:competitors} משווה את החלופות.

\begin{table}[htbp]
\centering\small
\rrstripes
\begin{tabularx}{\linewidth}{>{\bfseries}p{2.6cm}XXXXX}
\rrtablehead
\rrth{חלופה} & \rrth{טריות} & \rrth{מים} & \rrth{השקעה (\en{capex})} & \rrth{יכולת הרחבה} & \rrth{מחיר ללקוח} \\
גגון (גג עירוני) & קטיף ביום האספקה, מרחק מאות מטרים & מחזור סגור, חיסכון משמעותי מול שדה & בינונית: \RRcapexRoof{} אלפי \shekel{} לגג & גבוהה, שכפול גג אחר גג & פרמיה מתונה סביב \RRprice~\shekel{} לק"ג \\
חווה אנכית פנימית & גבוהה & מחזור סגור & גבוהה מאוד, כולל חשמל לתאורה & בינונית, דורשת נדל"ן ואנרגיה & גבוה \\
חממה מסחרית & בינונית, הובלה מהפריפריה & בינוני & בינונית & גבוהה & בינוני \\
תוצרת מיובאת או שדה & נמוכה עד בינונית & גבוה & נמוכה ללקוח & גבוהה מאוד & נמוך \\
ערכת בלקון ביתית & גבוהה, אך כמות זעירה & נמוך & נמוכה מאוד & אין & נמוך ליחידה, גבוה לק"ג \\
\end{tabularx}
\caption{השוואת גגון לקטגוריות חלופות (הערכה איכותית של הצוות).}
\label{tab:competitors}
\end{table}

היתרון של גגון איננו בעלות ללקוח אלא בשילוב של טריות, אספקה קצרה ושקיפות. מערכת הידרופונית \en{NFT} על גג נשענת על אור שמש לגידול, וחשמל נדרש בעיקר למשאבות ולחיישנים, ובמים ממוחזרים. הספרות על חקלאות אנכית וסביבה מבוקרת דנה בפוטנציאל ובמגבלות בתחום המים והאנרגיה \cite{benke2017}.

\section{ניתוח \en{SWOT}}\label{sec:swot}

איור~\ref{fig:swot} מסכם את החוזקות והחולשות הפנימיות ואת ההזדמנויות והאיומים החיצוניים. אנו רואים את הגג כנכס דל-עלות שאין לו שימוש אחר, ואת תוכנת \en{RootOS} כחפיר שמקשה על חיקוי. מנגד, התלות באישורי הבניין והעלות ההתחלתית לגג הם נקודות תורפה שהתוכנית מטפלת בהן.

\begin{figure}[htbp]
  \centering
  \resizebox{\linewidth}{!}{% SWOT 2x2 — internal (top) vs external (bottom), helpful (right, RTL) vs harmful (left).
\begin{tikzpicture}[
  q/.style={minimum width=7.1cm,minimum height=4.2cm,rounded corners=3mm,
            fill=#1,draw=none},
  t/.style={font=\sffamily\bfseries\Large,text=white,anchor=north east},
  b/.style={font=\small,text width=6.2cm,align=right,anchor=north east}]
  \node[q=rrForest]   (S) at (3.75,2.25) {};
  \node[q=rrTerra]    (W) at (-3.75,2.25) {};
  \node[q=rrLeaf!80!black] (O) at (3.75,-2.25) {};
  \node[q=rrStone]    (T) at (-3.75,-2.25) {};
  \foreach \n/\l in {S/חוזקות \en{S},W/חולשות \en{W},O/הזדמנויות \en{O},T/איומים \en{T}}
    \node[t] at ([xshift=-3mm,yshift=-2mm]\n.north east) {\l};
  \node[b,text=white] at ([xshift=-3mm,yshift=-11mm]S.north east)
   {$\bullet$\ מודל ללא קרקע — הגג כבר קיים\\ $\bullet$\ \en{RootOS} מקטין תלות במגדל מומחה\\ $\bullet$\ החזר השקעה לגג תוך \RRroofPayback{} שנים\\ $\bullet$\ פי 12 פחות מים מגידול בשדה};
  \node[b,text=white] at ([xshift=-3mm,yshift=-11mm]W.north east)
   {$\bullet$\ מותג חדש ללא היסטוריה\\ $\bullet$\ תלות באישור ועד בית\\ $\bullet$\ צריכת חשמל למשאבות ולחיישנים\\ $\bullet$\ תפוקה מוגבלת לגג בודד};
  \node[b,text=white] at ([xshift=-3mm,yshift=-11mm]O.north east)
   {$\bullet$\ תוכניות אקלים עירוניות וגגות ירוקים\\ $\bullet$\ ביקוש גובר לתוצרת מקומית וטרייה\\ $\bullet$\ מכירת \en{RootOS} לחוות אחרות\\ $\bullet$\ התרחבות לחיפה ולבאר שבע};
  \node[b,text=white] at ([xshift=-3mm,yshift=-11mm]T.north east)
   {$\bullet$\ גלי חום וקיץ ארוך\\ $\bullet$\ ירידת מחירי ירקות עלים\\ $\bullet$\ מתחרים עם הון גדול (חוות אנכיות)\\ $\bullet$\ שינויי רגולציה בבנייה};
\end{tikzpicture}
}
  \caption{ניתוח \en{SWOT} של גגון.}
  \label{fig:swot}
\end{figure}

\begin{riskbox}[סיכון תחרותי]
יזם חוות אנכית בנדל"ן זול או רשת קמעונאית גדולה יכולים להיכנס לנישת ה"מקומי והטרי". ההגנה של גגון: חוזים ל-\RRcontractYears{} שנים עם ועדי בתים, נתוני גידול של \en{RootOS} ויתרון מיקום ברדיוס הלקוח.
\end{riskbox}

\section{מה המשמעות למשקיע}

השוק מאפשר את התוכנית בלי הנחות שוק אגרסיביות. המתחרים החזקים יקרים להפעלה או רחוקים מהלקוח, ולכן גגון תופסת נישה ברורה. הצמיחה מוגבלת בעיקר על ידי קצב ההתקנה של הגגות, והפרקים הבאים עוסקים בו.
